\documentclass[12pt,a4paper]{article}

% --- pakiety językowe i kodowanie ---
\usepackage[utf8]{inputenc}
\usepackage[T1]{fontenc}
\usepackage[polish]{babel}
\usepackage{polski}

% --- pakiety graficzne i formatowanie ---
\usepackage{geometry}
\geometry{margin=2.5cm}
\usepackage{graphicx}
\usepackage{hyperref}
\usepackage{enumitem}
\usepackage{listings}
\usepackage{booktabs}
\usepackage{float}
% --- konfiguracja listingu (kodu) ---
\lstset{
    basicstyle=\ttfamily\small,
    breaklines=true,
    frame=single,
    captionpos=b
}

% --- dane o projekcie ---
\title{Dokumentacja Funkcjonalna Projektu C\\ \large System wyznaczania współrzędnych węzłów grafów planarnych}
\author{Autorzy: Mateusz Bombola, Kajetan Karwacki}
\date{\today}

\begin{document}

\maketitle
\newpage
\tableofcontents
\newpage

\section{Cel projektu}
Celem projektu jest stworzenie aplikacji konsolowej w języku C, która na podstawie topologii grafu planarnego (listy krawędzi) wyznacza optymalne współrzędne dwuwymiarowe $(x, y)$ dla jego wierzchołków. 

Głównym zadaniem programu jest przekształcenie danych o połączeniach (podanych w postaci danych tekstowych) w strukturę umożliwiającą wizualizację, która minimalizuje przecięcia krawędzi i zachowuje czytelność struktury grafu. Aplikacja stanowi silnik obliczeniowy, którego wyniki mogą być w dalszym etapie wykorzystane przez moduły graficzne.

\section{Opis funkcji obsługiwanych przez program}
Program posiada następujące funkcjonalności:
\begin{itemize}
    \item \textbf{Wczytywanie struktury grafu:} Przetwarzanie plików tekstowych zawierających listę krawędzi wraz z ich wagami.
    \item \textbf{Obliczanie układu grafu:} Wyznaczanie współrzędnych wierzchołków przy użyciu jednego z dwóch zaimplementowanych algorytmów:
    \begin{itemize}
        \item \textbf{Algorytm Fruchtermana-Reingolda:} Model siłowy (force-directed), traktujący krawędzie jak sprężyny, a wierzchołki jak ładunki odpychające się.
        \item \textbf{Algorytm Tutte’a (Spring Theorem):} Metoda gwarantująca poprawny układ dla grafów planarnych trójspójnych poprzez umieszczenie wierzchołków zewnętrznych na wielokącie wypukłym.
    \end{itemize}
    \item \textbf{Eksport wyników:} Zapis wyliczonych współrzędnych do pliku w formacie tekstowym lub binarnym.
\end{itemize}

\section{Interfejs linii poleceń}
Program jest sterowany wyłącznie poprzez argumenty przekazywane w wierszu poleceń. Składnia wywołania:
\begin{center}
    \texttt{./graph\_layout [FLAGI]}
\end{center}

\subsection{Lista dostępnych flag}
\begin{table}[H]
\centering
\begin{tabular}{@{}lll@{}}
\toprule
\textbf{Flaga} & \textbf{Argument} & \textbf{Opis} \\ \midrule
\texttt{-i} & \textit{ścieżka} & Ścieżka do wejściowego pliku tekstowego z grafem (wymagane). \\
\texttt{-o} & \textit{ścieżka} & Ścieżka do pliku wynikowego (wymagane). \\
\texttt{-a} & \textit{algorytm} & Wybór algorytmu: \texttt{fr} (Fruchterman-Reingold) lub \texttt{tutte}. \\
\texttt{-f} & \textit{format} & Format wyjściowy: \texttt{txt} (tekstowy) lub \texttt{bin} (binarny). \\
\texttt{-h} & --- & Wyświetlenie pomocy i zakończenie programu. \\
\bottomrule
\end{tabular}
\caption{Argumenty wywołania programu}
\end{table}

\section{Formaty danych}

\subsection{Format wejściowy (tekstowy)}
Program przyjmuje plik tekstowy, gdzie każda linia reprezentuje jedną krawędź w formacie:
\begin{lstlisting}
<nazwa_krawedzi> <wierzcholek_A> <wierzcholek_B> <waga_krawedzi>
\end{lstlisting}
\textbf{Przykład:}
\begin{lstlisting}
E1 1 2 1.0
E2 2 3 1.0
E3 3 1 1.5
\end{lstlisting}

\subsection{Format wyjściowy (tekstowy)}
Wynik działania zapisywany jest jako lista współrzędnych:
\begin{lstlisting}
<wierzcholek> <wspolrzedna_x> <wspolrzedna_y>
\end{lstlisting}
Współrzędne są liczbami zmiennoprzecinkowymi.

\section{Sytuacje wyjątkowe i błędy}
Program zwraca odpowiedni kod błędu (\texttt{exit code}) oraz komunikat na standardowe wyjście błędów (\texttt{stderr}).

\begin{itemize}
    \item \textbf{Kod 1: Błąd odczytu pliku} – Podana ścieżka wejściowa nie istnieje lub brak uprawnień.
    \item \textbf{Kod 2: Błędny format danych} – Plik wejściowy zawiera znaki niebędące liczbami tam, gdzie są wymagane, lub brakuje kolumn.
    \item \textbf{Kod 3: Nieznany algorytm} – Przekazano flagę \texttt{-a} z wartością inną niż \texttt{fr} lub \texttt{tutte}.
    \item \textbf{Kod 4: Błąd alokacji pamięci} – Brak pamięci RAM do przetworzenia dużego grafu.
\end{itemize}

\section{Przykłady użycia}
1. Wyznaczenie współrzędnych metodą Fruchtermana-Reingolda i zapis do tekstu:
\begin{lstlisting}
./graph_layout -i dane.txt -o wynik.txt -a fr -f txt
\end{lstlisting}
\noindent
2. Wyznaczenie współrzędnych metodą Tutte'a i zapis binarny (500 iteracji):
\begin{lstlisting}
./graph_layout -i siec.txt -o mapa.bin -a tutte -f bin -n 500
\end{lstlisting}

\end{document}